\chapter{57}

\section{Mailand (IT), Mittwoch, 22. 
August}

«Eine harte Nuss, arrogant und verschlossen, \emph{porca miseria}\footnote{«Heilige Scheisse.»}.»

Mantovani haute auf den Tisch. Sömnez fuhr zusammen. Rossi unterdrückte
ein Grinsen.
So kannten sie ihn nicht. Der Fall schien ihm zuzusetzen.

«Scusi.»

Mantovani blickte hoch und formte seine Hände zur Gebetshaltung. Er
liess die Daumen drehen.

«Ich wollte euch nicht erschrecken. Nur, es ist, wie ich befürchtete.
Der junge di Martino kreuzte gestern mit Vater und Anwalt auf. Es
handelt sich leider tatsächlich um den Sohn \emph{des} Bernardo di
Martino. Der Junge sprach kein Wort, schaute mich siegessicher an und
sein Anwalt bestritt für diesen besagten Abend einen näheren Kontakt mit
der Verunfallten. Und di Martino sei sogar vor allen andern gegangen, er
habe nur ein paar Worte mit der jungen Frau gewechselt. Obwohl, ihre
Beziehung sei gut und eng, eine Freundschaft seit Kindertagen. Dann
redete der Vater, kontrolliert, leise und messerscharf, bis ihn der
Anwalt beschwichtigte und das Wort wieder zurückholte.
Bastiano di Martino wirkte je länger je erstaunter, schien mir
jedenfalls, und als der Anwalt drohte, die Presse wegen Verleumdung
einzuschalten, juckte der Junge sichtlich hoch. Zum Schluss sagte der
Vater, ich wisse ganz genau, dass das alles Bullshit sei. Der Ausdruck
passte schlecht zu seinem schicken Outfit.»

Die beiden Ermittler grinsten. Mantovani atmete hörbar aus.

«Die sichergestellten Fingerabdrücke auf der Tasche und an den
gefundenen Kondomen beeindruckten die nicht. Da es bei Giorgia Sala keine
Anzeichen von Geschlechtsverkehr gab, hatte ich wenig
Spielraum. Aber, gegen eine Haaranalyse des Jungen in Bezug auf Drogen
hatten sie keine Chance. Unsere einzige Möglichkeit, etwas mehr zu
erfahren.»

Mantovani schien wieder der Alte zu sein.

Rossi holte Luft.

«Wie haben Sie die Analyse begründet?»

«Die Blutuntersuchungen von Giorgia Sala waren zum Glück schon im Laufe
des Nachmittags auf meinem Tisch. Zur Sicherheit läuft übrigens parallel
eine Haaranalyse. Die junge Frau hatte einen heftigen Drogencocktail
konsumiert. Ihre Eltern meinten, das erstaune sie. Aber ja, Eltern sind
immer mal wieder erstaunt über ihre erwachsenen Kinder. Wir werden uns
im nahen Umfeld der Verletzten umhören müssen, im Freundeskreis und so
weiter. Noch ist sie nicht ansprechbar. Den Staatsanwalt konnte ich
aufgrund des Drogennachweises leicht ins Boot holen. Er unterstützt uns.
Di Martinos Anwalt musste der Haaranalyse zustimmen. Der Vater vibrierte
vor Wut. Bis morgen früh haben wir die Auswertung.»

Sömnez machte ununterbrochen Notizen. Rossi stutzte.

«Drogen? Was heisst heftiger Drogencocktail?»

«Der höchste Anteil war PCP. Neben MDMA. PCP soll zu Leichtsinnigkeit
und Euphorie führen. Dem Hochgefühl folgen bei vielen Menschen
regelrechte Angst-Ausbrüche. Es kann zu verzerrter Wahrnehmung des
Körpers und der Umwelt kommen. Auch das Zeitgefühl ist weg.»

«Krass,» Sömnez zog die Nase kraus und kratzte sich an der Schläfe,
«wäre es möglich, dass di Martino ihr die Droge organisiert hat und
selbst womöglich auch was eingeworfen hatte?»

«Möglich, wobei wir damit nicht unbedingt weiter wären.»

Die Tür zu Mantovanis Büro wurde geöffnet und ein junger Mann erschien.

«Kaffee?»

Die drei schienen sich synchron zu entspannen.

«Grossartig, danke Aldo, du kommst wie gerufen.»

Der Mann stellte ein rundes Tablett mit drei Espressi zwischen die
Papiere auf den Tisch. Dann zog er die Mundwinkel hoch und verschwand
dezent wie er gekommen war. Mantovani und Sömnez tranken den Kaffee
schwarz. Rossi rührte gedankenversunken den Zucker ein.
Mantovani sammelte sich. 

«Nun zum unangenehmsten Teil, Ragazzi,» so nannte er seine Ermittler,
wenn er feststeckte, «Giorgia Sala ist in einem gefährlichen Zustand.
Davon habe ich natürlich in der Vernehmung gestern nichts durchsickern
lassen, im Gegenteil, ich habe geblufft. Di Martino ist im Glauben, dass
wir bereits mit der jungen Frau gesprochen haben.»

Er kratzte sich mit dem Zeigefinger im Haar. Sömnez und Rossi hingen ihm
wie Magnete an den Lippen.

«Heute morgen sprach ich mit dem behandelnden Arzt. Ihr Herz-
Kreislaufsystem wird nach wie vor durch intensivmedizinische Massnahmen
aufrecht erhalten. Sie ist bis auf weiteres im Koma, soweit nicht
ungewöhnlich nach einer Kopfverletzung dieser Art. Die letzte Nacht gab
es aber Komplikationen. Sie musste neu stabilisiert werden. Der Arzt
vermutet eine zusätzliche Verletzung des Rückenmarks.»

Mantovani schaute von seinen Notizen hoch und geriet mit seinem Ellbogen
an den kleinen Kaffeelöffel am Tischrand. Als wollte das Scheppern
seine Bekümmerung unterstreichen, war er auf den Steinboden
gefallen.

«Es sieht nicht gut aus.»

Rossi streckte sich wie immer, wenn seine Gedanken durcheinander kamen.

«Wie steht es um sie?»

«Hm, die Frau wird vorläufig nicht aufwachen und wenn die Vermutung des
Arztes sich als richtig erweist, folgt eine Rückenoperation. Ein
Spezialist wird morgen zugezogen. Spätestens morgen Abend sollten wir
mehr wissen.»

«Was heisst das genau?»

«Weiss ich nicht. Klar ist, dass wir sie noch immer nicht befragen
können. Wir brauchen dringend Zeugen. Wer war in der Nacht auf den 20.
August neben Bastiano di Martino am \emph{Roccia del Culto}?»

Rossi fuhr Mantovani ins Wort. Seine Ungeduld nahm Überhand.

«Was meinte di Martino dazu»?

«Natürlich nichts. Auch dazu schwieg er, leider,» Mantovani
räusperte sich, «bis jetzt hat sich auch auf unseren Zeugenaufruf
niemand gemeldet und über den Verrückten, der hier bei uns aufkreuzte,
wissen wir auch noch nichts. Die Videoüberwachung brachte nicht mehr
Hinweise, als uns Cristina bereits gab. In einem weiteren Schritt
bringen wir die Aufnahme an die Öffentlichkeit und machen einen Aufruf.
Doch wie immer, der letzte Schritt.»

Mantovanis Handy klingelte. Die Vibrationen liessen es auf dem Tisch
tanzen. Nach einem beiläufigen Blick auf das Display drückte er die
Nummer weg.

«Später, der Arzt war’s. Also, ihr beide besucht bitte Giorgia Salas
Eltern. Ruft vorher in der Tanzschule an. Versucht soviel wie möglich
über Freundinnen, Freunde und Bekannte zu erfahren. Stellt eine Liste
mit Namen und Handynummern zusammen. Nutzt den Tag. Befragt alle in
Bezug auf die Nacht am Felsen. Und holt Informationen über Bastiano di
Martino ein. Wenn er ein Freund aus Kindertagen ist, müssten Parallelen
zu andern Bekannten bestehen. Verfolgt jeden kleinsten Hinweis! Noch
Fragen?»
 
Sie hatten keine, standen auf und verabschiedeten sich.

Auf dem Vorplatz der \emph{Questur} wählte Sönmez die Nummer der \emph{Scuola di Danza}.
